% !TEX root = ../Tese_LucasMaio.tex

\chapter{Objetivos}

\section{Objetivo Geral}

Desenvolver uma interface gráfica interativa, em Python, voltada à reconstrução e visualização de dados de Tomografia por Impedância Elétrica (TIE), com foco educacional e de apoio à pesquisa, permitindo explorar comparativamente como diferentes métodos de reconstrução e configurações da malha afetam as imagens obtidas.

\section{Objetivos Específicos}

\begin{enumerate}

\item Desenvolver a aplicação EITduino, utilizando os métodos de reconstrução Back-Projection (BP), Jacobiano (JAC) e GREIT implementados pela biblioteca de código aberto \texttt{pyEIT}, com interface interativa para alternância entre métodos e visualização comparativa dos resultados.

\item Permitir a exploração interativa de parâmetros de reconstrução e de configuração da malha, como densidade, formato geométrico, resolução da grade de renderização e regularização, evidenciando o efeito dessas escolhas sobre as imagens obtidas.

\item Implementar uma segunda interface de visualização utilizando a biblioteca \texttt{PyQtGraph}, a fim de comparar qualitativamente essa abordagem com a interface principal baseada em \texttt{Matplotlib} em termos de organização visual e desempenho de renderização.

\item Disponibilizar publicamente o código-fonte da aplicação, organizado de forma modular e documentado de modo a facilitar sua reutilização, reprodução e extensão em trabalhos futuros.

\end{enumerate}